\section{Experimental Setup}
\label{sec:setup}

\subsection{Data and split}
\label{sec:data}

All arms train on WikiText-103 \citep{merity2017wikitext} with a byte-pair
vocabulary of 8{,}192 tokens and packed sequences of length 256
% >>>TMLRONLY
(\texttt{lang-\allowbreak wikitext-\allowbreak 103-\allowbreak bpe8192-\allowbreak len256-\allowbreak 800d6154}). We use the corpus's own
% <<<TMLRONLY
% >>>ACLONLY
% The literal identifier is 41 characters of \texttt, and TeX sets
% \hyphenchar to -1 for typewriter, so it cannot break anywhere and ran 98pt
% past the margin. Short phrase here, literal string in the appendix.
(identifier in Appendix~\ref{app:protocol}). We use the corpus's own
% <<<ACLONLY
author-provided validation split rather than carving one out of the training
material, which is what makes the leakage audit in \S\ref{sec:integrity}
meaningful: no validation or test token was ever available to the trainer.

The baseline floor is a backoff bigram estimated on the same corpus:
\textbf{4.9849 nats/token} (7.19 bits, perplexity 146.2). A model that does not
beat it has learned nothing a two-column count table could not. All three arms
beat it, by margins of $+1.037$, $+1.470$ and $+1.348$ nats for MoE, MoR and MoRE
respectively.

% >>>ACLMOVE:protocol|\subsection{Training protocol}\label{sec:protocol} Five seeds per arm; 3 epochs of AdamW (lr $10^{-3}$, weight decay $10^{-4}$, clip 1.0, dropout 0.1), $d_{\text{model}}=256$, one block, 4 heads, $T=7$, batch 48, each run on a single GPU (\S\ref{sec:honesty}). Settings were fixed once, before the comparison, and not swept per architecture. \textbf{Three epochs is a fixed compute budget, not convergence}---every arm gets the identical budget, so the comparison is valid though no arm is converged (Appendix~\ref{app:protocol}).
\subsection{Training protocol}
\label{sec:protocol}

% >>>ACLONLY
% The ACL body replaced the 41-character dataset identifier with a pointer here,
% because \texttt cannot hyphenate and the literal string ran 98pt past the
% margin. It must therefore actually appear in the appendix. Nested inside the
% ACLMOVE:protocol block on purpose: the assembler resolves variants BEFORE
% extracting moves, so this lands in the ACL appendix and is absent from TMLR,
% which states the identifier in its own Data section.
The frozen dataset identifier is
\texttt{lang-\allowbreak wikitext-\allowbreak 103-\allowbreak bpe8192-\allowbreak len256-\allowbreak 800d6154}, and the split is the
corpus's own author-provided validation split.
% <<<ACLONLY

Every canonical arm runs at seeds $\{42,43,44,45,46\}$ and is reported as
mean$\,\pm\,$std over five runs. Training is 3 epochs of AdamW (lr $10^{-3}$,
weight decay $10^{-4}$, gradient clipping 1.0, dropout 0.1) with
$d_{\text{model}} = 256$, one block, 4 attention heads, $T = 7$ and batch size 48,
each run on a single GPU: an RTX 4060 laptop GPU for the two arms retrained after
a checkpoint loss, a V100 for the composed arm (\S\ref{sec:honesty}). These
settings were fixed once, before the comparison, and
not swept per architecture: they are a stable operating point, and no arm was tuned
for itself.

\textbf{Three epochs is a fixed compute budget, not convergence, and we state it
as such.} One epoch is 526{,}320 blocks $\times$ 256 tokens $= 134.7$M tokens
against $\approx$3.4M non-embedding parameters, so the runs sit in an underfitting
regime; the choice of 3 over 2 was made because the halting axis needs more than
two epochs to separate the arms at all, and going to 8 epochs was rejected on cost.
Every arm receives the identical budget, so the comparison is valid even though no
arm is converged. The consequences are stated in \S\ref{sec:limitations}.

% <<<ACLMOVE:protocol
% >>>ACLMOVE:metrics|\subsection{Metrics}\label{sec:metrics} The primary metric is \textbf{validation task loss}, final-epoch per-token cross-entropy in nats; \texttt{best} equals \texttt{last} on all 15 cells, so \emph{no model selection occurred}. Routing is scored with permutation-invariant measures (AMI, Hungarian, purity) against a per-run shuffled-label control and reported as a z-score. Depth is reported with its exit rates alongside: a depth--difficulty correlation on an arm whose tokens almost all reach the ceiling is taken against a near-constant and \textbf{says nothing about allocation}. Compute is inference FLOPs per token, counted on a real forward pass at each arm's measured mean depth (Appendix~\ref{app:flops}). Appendix~\ref{app:metrics} defines the rest.
\subsection{Metrics}
\label{sec:metrics}

\textbf{Primary metric: validation task loss}, the final-epoch per-token
cross-entropy in nats. We deliberately do not select on best-epoch validation loss:
that is an order statistic whose downward bias scales with each arm's per-epoch
validation noise, which is not matched across architectures. On all 15 cells
\texttt{best} equals \texttt{last}---the evaluation interval is one epoch and the
curves decrease monotonically---so \emph{no model selection occurred} on either
split. Best-epoch values are reported alongside so the selection effect stays
checkable.

Because the router is unsupervised, expert indices carry no semantic identity, so
raw agreement with any externally defined partition measures label permutation as
much as partition quality. We therefore report AMI \citep{vinh2010ami},
Hungarian-matched accuracy and cluster purity, each against a \emph{per-run
permutation control}: the same metric computed against shuffled linguistic labels,
averaged over 10 draws. The controlled quantity is the z-score of the real value
against that null, not the raw value. The externally defined partition is a
part-of-speech hypothesis about a useful expert split---a hypothesis, not a
functional ground truth---so a router can be better for next-token prediction and
still score low against it. We therefore never call this ``routing accuracy''
without that qualifier.

For depth we report mean steps, exit rates, halt mass, and the Spearman correlation
between per-token depth and per-token difficulty. \textbf{The correlation is only
interpretable when depth has variance to correlate}: on an arm whose tokens almost
all exit at the ceiling, the correlation is computed against a near-constant and a
value near zero is uninformative rather than negative evidence. We report the exit
rates next to every correlation for that reason (\S\ref{sec:halting}).

Compute is reported as inference FLOPs per token, counted on a real forward pass
and evaluated at each arm's measured mean depth; Appendix~\ref{app:flops} gives the
method and its checks. FLOPs are an analytic cost, not a runtime, so measured
throughput is reported beside them.

% <<<ACLMOVE:metrics
% >>>ACLMOVE:flops|
We count FLOPs rather than derive them. \texttt{torch.utils.flop\_counter} records
every operation that executes on a real forward pass of each arm, built from its
own resolved configuration, so the count cannot drift from the model: attention,
the embedding and positional tables, the router, the per-expert halt heads and the
8{,}192-way output layer are all included, and none is approximated by a formula.
Depth is swept by forcing every token to exactly $k$ steps for $k = 1,\dots,7$.
The recursive part of the graph is linear in $k$, so
$\mathrm{FLOPs}(d) = \mathrm{FLOPs}(1) + (d-1)\,\Delta$, with $\Delta$ the marginal
cost of one step; the fit is \emph{verified} against the measured $k = 7$ count, and
a residual above $0.1\%$ is treated as a failure, because a non-linearity there
would mean the recursion is not computing what the architecture claims. The
residual is zero for both recursive arms.

Each arm is then costed at its \emph{measured} mean depth---the mean over its five
admitted runs of the realised number of steps per token (MoR $6.730$, MoRE
$6.955$)---not at the depth budget, which would overstate both recursive arms by
whatever early exit saves. Because FLOPs are linear in depth, the cost at the mean
depth equals the mean of the per-run costs. Per token this gives
$5.25$\,M (MoE), $46.5$\,M (MoR) and $11.5$\,M (MoRE): MoR costs $8.87\times$ MoE
and MoRE $2.19\times$, so MoR costs $4.05\times$ MoRE. These are the FLOPs a
halting-aware implementation would execute. Ours executes somewhat more for both
recursive arms: the expert FFN runs only for still-active tokens, but attention is
computed at every position for as long as any token in the batch is active, and
its output is discarded at halted ones---a cost the engine logs rather than hides.
Both recursive arms run most batches to the full budget, so the excess falls on
both, and slightly more on MoR, whose mean depth is lower.

Throughput is measured separately, because FLOPs are not a runtime and the two
disagree here. The admitted runs trained on two different GPUs, so their logged
training throughputs compare hardware as much as architectures; we therefore time
inference for all fifteen checkpoints on one GPU (an RTX 3050 laptop GPU), on the
identical out-of-domain blocks at batch size 32, after a warm-up batch.
% <<<ACLMOVE:flops
\subsection{Statistics}
\label{sec:stats}

Comparisons use an exact two-sided randomization test on the difference of means,
which enumerates every way of splitting the pooled values into groups of the
original sizes---$\binom{10}{5} = 252$ at five-versus-five---and counts the
fraction whose $|$mean difference$|$ is at least the observed one. It is
exhaustive, hence deterministic: no sampling error, no seed, no normality
assumption, and no threshold to choose. The superseded $k\times$std heuristic is
not used anywhere.

The \textbf{resolution floor} of this design --- the smallest $p$ it can return at
all --- is $2/\binom{n_a+n_b}{n_a}$ for equal arms, which is $2/252 = 0.00794$ at
five-versus-five. The factor of two is there because the two-sided statistic is
invariant under swapping the groups, so extreme regroupings arrive in mirror pairs
and the count at the extreme can never be one.
% >>>ACLMOVE:floor_proof|The derivation is in Appendix~\ref{app:floor}.
The argument is worth stating because the floor is a property of the design rather
than of the data. Our statistic is $|\mathrm{mean}(a) - \mathrm{mean}(b)|$. When
the two arms are the same size, the complement of any admissible group-$a$ subset
is itself an admissible group-$a$ subset, and swapping the groups negates the
difference while preserving its absolute value. Every regrouping at the observed
extreme therefore has a partner that is equally extreme, the count can never be
odd at the tail, and the smallest attainable fraction is $2/\binom{n_a+n_b}{n_a}$
rather than $1/\binom{n_a+n_b}{n_a}$. When the arms differ in size the complement
has the wrong size, is not an admissible regrouping, no mirror is forced, and
$1/\binom{n_a+n_b}{n_a}$ becomes attainable.
% <<<ACLMOVE:floor_proof
\textbf{Every headline comparison in this paper sits exactly on that floor}, which
we phrase as ``as extreme as a five-seed design admits''---a weaker statement than
significance, and not to be dressed up as a stronger one. Only more seeds can
strengthen these claims; no larger effect can.

% >>>ACLMOVE:holm_detail|The correction is Holm--Bonferroni over two pre-registered families (predictive quality, $m=4$; adaptive depth, $m=2$), declared before the final arm was trained. Every headline comparison clears its threshold, and would clear plain Bonferroni even with all six declared comparisons pooled into one family ($0.05/6=0.00833$, a margin of 0.0004 over the floor); Appendix~\ref{app:stats} gives the memberships.
Correction is Holm--Bonferroni over two \emph{pre-registered} families, declared
before the final arm was trained: Family A (predictive quality, $m=4$, metric
validation task loss) and Family B (adaptive depth, $m=2$, metric depth--difficulty
correlation). The split is declared rather than chosen because the two families ask
different questions on different quantities, and pooling them would correct the
depth claim for tests that cannot bear on it. Holm controls the same familywise
error rate as plain Bonferroni and is uniformly more powerful, but at our
resolution floor the two give identical verdicts: Holm's first-step threshold at a
given $m$ \emph{is} Bonferroni's. What sets the threshold is the family size,
$0.05/4 = 0.0125$ in Family A and $0.025$ in Family B. Even pooling all six declared
comparisons into one family gives $0.05/6 = 0.00833$, which the floor of $0.00794$
still clears---by $0.0004$. That margin is the binding constraint on this design: a
seventh confirmatory comparison would put the floor above the threshold, after
which no effect of any size could be declared significant. It is one more reason
no comparison may be added to a family once results exist.

Two consequences we hold ourselves to. A declared comparison that could not be run
still consumes its slot: dropping a failed arm must never make the survivors easier
to call significant. And a comparison whose floor exceeds its Holm threshold is
\emph{undecidable} and reported \texttt{N/A}, never ``not significant''---that
phrasing would blame the data for a limit of the design.
% <<<ACLMOVE:holm_detail

% >>>ACLMOVE:integrity|\subsection{Integrity checks}\label{sec:integrity} Three invariants hold across the matrix, each enforced by an automated check: inputs are \textbf{bit-identical across arms} for the same record; the target is absent from the input and no oracle routing label is encoded in a feature slot; and the arms are parameter-matched to $0.069\%$ with seed-blind configurations identical within each arm, so comparability rests on a frozen specification rather than on one code revision. Appendix~\ref{app:integrity} lists them with the structurally-zero loss terms.
\subsection{Integrity checks}
\label{sec:integrity}

Three mechanical invariants hold across the matrix, each enforced by an automated
check rather than by inspection. \textbf{(i) Inputs are bit-identical across arms
for the same record}: feature ranges are identical at $E \in \{1,5,6\}$, so a
difference between arms cannot be a difference in representation. \textbf{(ii) The
regression or classification target is absent from the input}, and no oracle routing
label is encoded into a feature slot. \textbf{(iii) The arms are parameter-matched
to 0.069\%} and the seed-blind configuration is identical within each arm, so
comparability cannot rest on a single code revision: the matrix spans four commits,
and it is the frozen specification plus the config-equality check that holds the
arms comparable, not one revision.

One disclosure belongs with the config table. Three loss terms are structurally
zero on language: the family-classification term because its head reads a detached
representation so its gradient cannot reach the trunk, and the routing-supervision
and halting-supervision terms because natural language supplies no per-token oracle
expert and no per-token oracle depth. The first is a genuine advantage over a
comparable synthetic study, where a family cross-entropy is often one of the larger
terms in the objective and every specialisation claim must be caveated for it.
% <<<ACLMOVE:integrity
